% Options for packages loaded elsewhere
\PassOptionsToPackage{unicode}{hyperref}
\PassOptionsToPackage{hyphens}{url}
\documentclass[
]{article}
\usepackage{xcolor}
\usepackage[margin=1in]{geometry}
\usepackage{amsmath,amssymb}
\setcounter{secnumdepth}{-\maxdimen} % remove section numbering
\usepackage{iftex}
\ifPDFTeX
  \usepackage[T1]{fontenc}
  \usepackage[utf8]{inputenc}
  \usepackage{textcomp} % provide euro and other symbols
\else % if luatex or xetex
  \usepackage{unicode-math} % this also loads fontspec
  \defaultfontfeatures{Scale=MatchLowercase}
  \defaultfontfeatures[\rmfamily]{Ligatures=TeX,Scale=1}
\fi
\usepackage{lmodern}
\ifPDFTeX\else
  % xetex/luatex font selection
\fi
% Use upquote if available, for straight quotes in verbatim environments
\IfFileExists{upquote.sty}{\usepackage{upquote}}{}
\IfFileExists{microtype.sty}{% use microtype if available
  \usepackage[]{microtype}
  \UseMicrotypeSet[protrusion]{basicmath} % disable protrusion for tt fonts
}{}
\makeatletter
\@ifundefined{KOMAClassName}{% if non-KOMA class
  \IfFileExists{parskip.sty}{%
    \usepackage{parskip}
  }{% else
    \setlength{\parindent}{0pt}
    \setlength{\parskip}{6pt plus 2pt minus 1pt}}
}{% if KOMA class
  \KOMAoptions{parskip=half}}
\makeatother
\usepackage{color}
\usepackage{fancyvrb}
\newcommand{\VerbBar}{|}
\newcommand{\VERB}{\Verb[commandchars=\\\{\}]}
\DefineVerbatimEnvironment{Highlighting}{Verbatim}{commandchars=\\\{\}}
% Add ',fontsize=\small' for more characters per line
\usepackage{framed}
\definecolor{shadecolor}{RGB}{248,248,248}
\newenvironment{Shaded}{\begin{snugshade}}{\end{snugshade}}
\newcommand{\AlertTok}[1]{\textcolor[rgb]{0.94,0.16,0.16}{#1}}
\newcommand{\AnnotationTok}[1]{\textcolor[rgb]{0.56,0.35,0.01}{\textbf{\textit{#1}}}}
\newcommand{\AttributeTok}[1]{\textcolor[rgb]{0.13,0.29,0.53}{#1}}
\newcommand{\BaseNTok}[1]{\textcolor[rgb]{0.00,0.00,0.81}{#1}}
\newcommand{\BuiltInTok}[1]{#1}
\newcommand{\CharTok}[1]{\textcolor[rgb]{0.31,0.60,0.02}{#1}}
\newcommand{\CommentTok}[1]{\textcolor[rgb]{0.56,0.35,0.01}{\textit{#1}}}
\newcommand{\CommentVarTok}[1]{\textcolor[rgb]{0.56,0.35,0.01}{\textbf{\textit{#1}}}}
\newcommand{\ConstantTok}[1]{\textcolor[rgb]{0.56,0.35,0.01}{#1}}
\newcommand{\ControlFlowTok}[1]{\textcolor[rgb]{0.13,0.29,0.53}{\textbf{#1}}}
\newcommand{\DataTypeTok}[1]{\textcolor[rgb]{0.13,0.29,0.53}{#1}}
\newcommand{\DecValTok}[1]{\textcolor[rgb]{0.00,0.00,0.81}{#1}}
\newcommand{\DocumentationTok}[1]{\textcolor[rgb]{0.56,0.35,0.01}{\textbf{\textit{#1}}}}
\newcommand{\ErrorTok}[1]{\textcolor[rgb]{0.64,0.00,0.00}{\textbf{#1}}}
\newcommand{\ExtensionTok}[1]{#1}
\newcommand{\FloatTok}[1]{\textcolor[rgb]{0.00,0.00,0.81}{#1}}
\newcommand{\FunctionTok}[1]{\textcolor[rgb]{0.13,0.29,0.53}{\textbf{#1}}}
\newcommand{\ImportTok}[1]{#1}
\newcommand{\InformationTok}[1]{\textcolor[rgb]{0.56,0.35,0.01}{\textbf{\textit{#1}}}}
\newcommand{\KeywordTok}[1]{\textcolor[rgb]{0.13,0.29,0.53}{\textbf{#1}}}
\newcommand{\NormalTok}[1]{#1}
\newcommand{\OperatorTok}[1]{\textcolor[rgb]{0.81,0.36,0.00}{\textbf{#1}}}
\newcommand{\OtherTok}[1]{\textcolor[rgb]{0.56,0.35,0.01}{#1}}
\newcommand{\PreprocessorTok}[1]{\textcolor[rgb]{0.56,0.35,0.01}{\textit{#1}}}
\newcommand{\RegionMarkerTok}[1]{#1}
\newcommand{\SpecialCharTok}[1]{\textcolor[rgb]{0.81,0.36,0.00}{\textbf{#1}}}
\newcommand{\SpecialStringTok}[1]{\textcolor[rgb]{0.31,0.60,0.02}{#1}}
\newcommand{\StringTok}[1]{\textcolor[rgb]{0.31,0.60,0.02}{#1}}
\newcommand{\VariableTok}[1]{\textcolor[rgb]{0.00,0.00,0.00}{#1}}
\newcommand{\VerbatimStringTok}[1]{\textcolor[rgb]{0.31,0.60,0.02}{#1}}
\newcommand{\WarningTok}[1]{\textcolor[rgb]{0.56,0.35,0.01}{\textbf{\textit{#1}}}}
\usepackage{graphicx}
\makeatletter
\newsavebox\pandoc@box
\newcommand*\pandocbounded[1]{% scales image to fit in text height/width
  \sbox\pandoc@box{#1}%
  \Gscale@div\@tempa{\textheight}{\dimexpr\ht\pandoc@box+\dp\pandoc@box\relax}%
  \Gscale@div\@tempb{\linewidth}{\wd\pandoc@box}%
  \ifdim\@tempb\p@<\@tempa\p@\let\@tempa\@tempb\fi% select the smaller of both
  \ifdim\@tempa\p@<\p@\scalebox{\@tempa}{\usebox\pandoc@box}%
  \else\usebox{\pandoc@box}%
  \fi%
}
% Set default figure placement to htbp
\def\fps@figure{htbp}
\makeatother
\setlength{\emergencystretch}{3em} % prevent overfull lines
\providecommand{\tightlist}{%
  \setlength{\itemsep}{0pt}\setlength{\parskip}{0pt}}
\usepackage{bookmark}
\IfFileExists{xurl.sty}{\usepackage{xurl}}{} % add URL line breaks if available
\urlstyle{same}
% fallback for those not using the hyperref driver hyperxmp:
\makeatletter
\@ifundefined{xmpquote}{\newcommand{\xmpquote}[1]{#1}}{}
\makeatother
\hypersetup{
  pdftitle={CS5702 Week1 Lab Worksheet},
  pdfauthor={Martin Shepperd},
  hidelinks,
  pdfcreator={LaTeX via pandoc}}

\title{CS5702 Week1 Lab Worksheet}
\author{Martin Shepperd}
\date{10/09/2026}

\begin{document}
\maketitle

This Worksheet is organised as an RMarkdown document. You can
\textbf{read} it. You can \textbf{run} the embedded R and you can
\textbf{add} your own R. We suggest you save it as another file so, if
necessary, you can revert to the original.

Whenever you click \textbf{Knit} in the RStudio menu it will render the
markdown into nicely formatted html (or pdf or even docx) which you can
examine in the Viewing Window in RStudio or via any web browser. You may
find this easier to read, however, you must edit the .rmd file, i.e.,
the RMarkdown in the Edit Pane if you want to make any changes. You are
encouraged to \emph{explore} and \emph{experiment.} Change something and
see what happens!

\textbf{Note}, that many of the topics and challenges are the same or
closely interwoven with CS5701 (QDA) so never be afraid to carry over
lessons and ideas from one module to the other.

\subsection{Pre-requisites}\label{pre-requisites}

\begin{enumerate}
\def\labelenumi{\arabic{enumi}.}
\tightlist
\item
  You should have completed the Week 0
  \href{https://martinshepperd-moderndatabook.share.connect.posit.cloud/00-getting-ready.html}{``Getting
  Ready''} and minimally must
  \href{https://martinshepperd-moderndatabook.share.connect.posit.cloud/00-getting-ready.html\#sec-C0_SoftwareInstall}{Install
  and test R and RStudio} material.
\item
  You should have attended/read the Week 1 Lectures.
\end{enumerate}

\section{0. Worksheet Topics}\label{worksheet-topics}

\begin{enumerate}
\def\labelenumi{\arabic{enumi}.}
\tightlist
\item
  \hyperref[WS1]{Create, save and edit R Notebooks using RStudio}
\item
  \hyperref[WS2]{Exploring the use of R using R}
\item
  \hyperref[WS3]{Some Basic R practice}\\
\end{enumerate}

\begin{enumerate}
\def\labelenumi{\Alph{enumi}.}
\tightlist
\item
  \hyperref[YAML_appendix]{\emph{Extension}: Appendix - notes on YAML}\\
\item
  \hyperref[ggplot_appendix]{\emph{Extension}: Appendix - notes on
  \{ggplot2\}}
\end{enumerate}

\subsection{0.1 Worksheet Navigation: Choose your track before
starting}\label{worksheet-navigation-choose-your-track-before-starting}

Welcome! Depending on your prior R experience, follow the appropriate
pathway below:

\textbf{TRACK A: New to R / Refreshing Core Concepts}\\
\textbf{Focus:} Work step by step through \textbf{Sections 1, 2 and 3}.
\textbf{Goal:} Take your time with the RStudio mechanics, running
RMarkdown code chunks, basic variable assignment and fundamental
functions. Do not feel pressured to reach the extension tasks (today).

\textbf{TRACK B: Experienced R Users (e.g., January Starters )}
\textbf{Focus:} Fast-track through \textbf{Sections 1 and 3} but work
through all of Section 2 as it's based on an actual data analysis
example. Target all \textbf{Extension Exercises} (e.g., \emph{Exercise
2.6} and \emph{Appendices 1 and 2} for YAML and \texttt{\{ggplot2\}}
visualisation work).\\
\textbf{Goal:} Focus on critical evaluation (e.g., \emph{Exercises 2.1
\& 2.3}) and converting base R graphics into something more professional
with \{ggplot2\}. Over the course of the next couple of weeks we will
all fully transition to \{ggplot2\}.

\emph{Note:} Both tracks use this single \texttt{.Rmd} file. Track B
extends the core tasks rather than skipping key learning objectives.

\section{1. Lab: create, save and edit RMarkdown documents using
RStudio}\label{WS1}

This worksheet is written using RMarkdown and saved as an
\href{https://blog.rstudio.com/2016/10/05/r-notebooks/}{RMD file}.

This is an extremely convenient way to combine explanatory text written
in a very simple markdown language and interactive R. We will use
RMarkdown regularly for all non-trivial R code in both modules (CS5701
and CS5702), including for labs and also for your coursework (CS5801)
and exam (CS5802). So it's important you become comfortable with using
RMarkdown documents.

If you've not already viewed my video ``Introducing RStudio'' or you
lack confidence in using this important tool I suggest you should check
it
\href{https://www.loom.com/share/f0ecbbbc17dc43acbcdf9deb5fb22f3c}{here}.

\begin{enumerate}
\def\labelenumi{\arabic{enumi}.}
\tightlist
\item
  You need a YAML header (see this file for an example, or click
  \hyperref[YAML_appendix]{here} for more information).\\
\item
  You can enter plain text as-is. Any simple editor will do or write
  them directly in RStudio. \emph{Avoid} word processors such as Word
  because the text is not stored as plain text and includes hidden
  non-printing characters which means you will get lots of truly
  horrible side effects!\\
\item
  You can use \texttt{\#}, \texttt{\#\#}, etc for different heading
  levels. Notice how RStudio automatically generates a nested Outline
  (navigation) Pane for you.\\
\item
  If you start a line with \texttt{1.} etc this generates an enumerated
  list (as per this example). There must be a space between the full
  stop and the text. If you just want bullets you can use \texttt{-}
  instead.

  \begin{itemize}
  \tightlist
  \item
    You can nest lists by indenting by 4 spaces.\\
  \end{itemize}
\item
  You can insert a code chunk by clicking the green \emph{Insert (Code
  Chunk)} button on the Edit Pane toolbar or by pressing
  \emph{Cmd+Option+I}. Or you could, if you really want to do it
  yourself, type the three backticks followed by the language in braces,
  then the code and close the chunk with another three backticks.\\
\item
  You can also view and edit your RMarkdown in `Visual' mode which shows
  the result of your formatting. It's your choice but I suggest you
  stick with `Source' mode to start with until you are confident editing
  the RMarkdown.
\end{enumerate}

There's also a very helpful guide on RMarkdown from RStudio
\href{https://rmarkdown.rstudio.com/lesson-2.html}{here}.

Below is an example of an R code chunk which can be executed from within
the Rmd document. You can do this by clicking the green triangle to the
top right of the chunk area within the document. If you wanted to run
all the code up until, and including, this code chunk you can click the
grey triangle with the green bar \footnote{In this case there is no
  previous code so it would make no difference but often RMarkdown
  contains multiple code chunks.}.

\begin{Shaded}
\begin{Highlighting}[]
\CommentTok{\# This is an R code chunk}
\FunctionTok{print}\NormalTok{(}\StringTok{"Goodbye cruel world"}\NormalTok{)}
\end{Highlighting}
\end{Shaded}

\begin{verbatim}
## [1] "Goodbye cruel world"
\end{verbatim}

\textbf{Exercise 1.1:} Edit the R code to output ``Goodbye cruel world''
or some other message. Re-run the code chunk.

\textbf{Exercise 1.2:} Save the edited RMarkdown document When you click
the \emph{Knit} button you can generate HTML which can be previewed in
the Viewer pane - bottom right). The raw R markdown has a suffix .Rmd
and the rendered HTML will have a suffix .html. Note that you can Knit
to other formats if you prefer, such as pdf \footnote{The way the file
  is rendered into pdf relies upon you having an installation of LaTeX
  on your machine. This tends to stress our lab machines - sorry! So if
  this doesn't work for you just focus on the html or reading the
  RMarkdown directly.}.

Check you can see the html version. The preview shows you a rendered
HTML copy of the contents of the editor. The output from the R code
chunks is whatever was generated when it was last run in the RStudio.
You can also view the file in any web browser.

\textbf{Exercise 1.3:} Add a new 3rd level heading named ``My RMarkdown
Exercise'' here. Add an anchor label
\texttt{\{\#\textless{}anchor-name\textgreater{}\}}. Now at another
point in the document you can link to your anchor, as follows:
\texttt{{[}text\ to\ display{]}(\#anchor-name)}. Lines 24-28 of this
RMarkdown file are examples.

\textbf{Exercise 1.4:} Create a new RMarkdown file using the
\texttt{File\ \textgreater{}\ New\ File\ \textgreater{}\ R\ Notebook}
menu option. RStudio provides you with a basic template which you can
modify as needed. Add a meaningful title and your name as an author in
the YAML header. Knit the new file using the \texttt{Knit} pull down
menu.

Click on my
\href{https://www.loom.com/share/13100c7fb0ef4e9cb453e1ed12ac73c8}{video
here} if you want further guidance on creating and manipulating
RMarkdown documents.

\textbf{Extension exercise 1.5:} You can also create a special type of
Rmd document called an R Notebook (analogous to a Jupityer notebook for
python). The main syntactic difference is in the choice of output in the
YAML header so you will have \texttt{output:\ html\_notebook}. The
advantage of notebook files is that code chunks can be run interactively
rather than as a batch when rendering (knitting) in Rmd documents. To
explore the differences further, see the tutorial from posit
\href{https://rmarkdown.rstudio.com/lesson-10.html}{here}.

To dig deep there's a nice blog of
\href{https://appsilon.com/r-markdown-tips/}{R Markdown Tips}.

So now you can handle the RMarkdown basics. Let's move onto R.

\section{2. Exploring the use of R using R}\label{WS2}

\subsection{2.1 Demand for data science software
skills}\label{demand-for-data-science-software-skills}

How useful are R skills and tools? We will use a recent analysis of US
job adverts in 2025. Source:
\href{https://www.onetonline.org/link/demand/15-2051.00}{Lightcast}
external site job postings data for the US between January 1, 2025 and
December 31, 2025. The percentage is the ratio of unique postings which
mention the skill to all unique postings linked to the 0*NET-SOC
occupation (categorised as data science).

The top 10 tools/technologies required by the employers are:

\begin{verbatim}
Python (66%)
SQL (51%)
R (34%)
Tableau (22%)
Microsoft Power BI (19%)
Amazon Web Services (17%)
Microsoft Azure (13%)
TensorFlow (11%)
PyTorch (10%)
SAS (9%)
\end{verbatim}

\textbf{Exercise 2.1:} It's important to look at the data critically,
though we must accept nothing is certain or perfect! What questions or
concerns might you have?

Since we're thinking about data and data analysis, it would be helpful
to visualise the data. Therefore I have copied some of the data from
their article (you can download a csv as well if you want) and plugged
this into R (see below). Don't worry about the coding details for the
present although be aware that for reasons of simplicity I'm just using
Base R for the graphics (functions like \texttt{pie()} and
\texttt{barplot()}). Later we will introduce the \{ggplot2\}\footnote{I,
  and many authors, use the convention of curly braces to denote an R
  package, so \{ggplot2\} means the package named ggplot2 which you
  could install with \texttt{install.packages("ggplot2")} and then read
  into memory with \texttt{library(ggplot2)}. If you're already an
  expert you can look at Appendix 2 to see how to do a pie chart with
  \{ggplot2\}, however, be warned it's a bit fiddly and also you're not
  really encouraged to make pie charts due to their many disadvantages.
  Barplots might be a better choice!} package which will enable almost
complete control and flexibility in producing beautiful data
visualisations. However, if you are already a confident R programmer
feel free to jump ahead.

If you want to know more about installing or using packages in the R
ecosystem see my
\href{https://www.loom.com/share/87b088c301cd40ac9aa4ff714141a6f3}{video
here}.

\begin{Shaded}
\begin{Highlighting}[]
\CommentTok{\# Reproduce the 2025 data from posted by 0*Net Online  }
\CommentTok{\# data on required tools for Data Scientists}
\CommentTok{\# supplied by Lightcast.}

\NormalTok{skill }\OtherTok{\textless{}{-}} \FunctionTok{c}\NormalTok{(}
  \StringTok{"Python"}\NormalTok{, }\StringTok{"SQL"}\NormalTok{, }\StringTok{"R"}\NormalTok{, }\StringTok{"Tableau"}\NormalTok{, }\StringTok{"Power BI"}\NormalTok{,}
  \StringTok{"AWS"}\NormalTok{, }\StringTok{"Azure"}\NormalTok{, }\StringTok{"TensorFlow"}\NormalTok{, }\StringTok{"PyTorch"}\NormalTok{, }\StringTok{"SAS"}
\NormalTok{)}
\NormalTok{percent }\OtherTok{\textless{}{-}} \FunctionTok{c}\NormalTok{(}\DecValTok{66}\NormalTok{, }\DecValTok{51}\NormalTok{, }\DecValTok{34}\NormalTok{, }\DecValTok{22}\NormalTok{, }\DecValTok{19}\NormalTok{, }\DecValTok{17}\NormalTok{, }\DecValTok{13}\NormalTok{, }\DecValTok{11}\NormalTok{, }\DecValTok{10}\NormalTok{, }\DecValTok{9}\NormalTok{)}
\CommentTok{\# Generate the pie chart}
\FunctionTok{pie}\NormalTok{(percent, }\AttributeTok{main=}\StringTok{"Pie chart of in demand data scientist tools"}\NormalTok{, }\AttributeTok{labels=}\NormalTok{skill)}
\end{Highlighting}
\end{Shaded}

\pandocbounded{\includegraphics[keepaspectratio]{CS5702_W1_Lab_files/figure-latex/unnamed-chunk-2-1.pdf}}

\textbf{Exercise 2.2:} Using the \texttt{clockwise} parameter change the
order of display. Type \texttt{?pie} for help if required.

\textbf{Exercise 2.3:} If you look at the help for the \texttt{pie()}
function you will read ``Pie charts are a very bad way of displaying
information.'' Why do you think this is so? {[}Hint: what might be the
impact of multi-valued data, e.g., many adverts will ask for more than
one skill? Also, how easy is it to read the values?{]}

\begin{Shaded}
\begin{Highlighting}[]
\CommentTok{\# NB we have already reproduced the data: percent and skill come from the }
\CommentTok{\# previous piechart chunk}

\CommentTok{\# Generate a bar chart}
\FunctionTok{barplot}\NormalTok{(percent, }
\AttributeTok{main=}\StringTok{"Bar chart of in{-}demand Data Scientist tools (2025)"}\NormalTok{, }
\AttributeTok{sub =} \StringTok{"Source: https://www.onetonline.org/link/demand/15{-}2051.00"}\NormalTok{,}
\AttributeTok{ylab=}\StringTok{"Percentage of listings that mention tool"}\NormalTok{, }
\AttributeTok{names.arg=}\NormalTok{skill, }
\AttributeTok{cex.names =} \FloatTok{0.6}\NormalTok{,       }\CommentTok{\# shrinks text size so fits on x{-}axis}
\AttributeTok{col=}\StringTok{"darkolivegreen"}\NormalTok{,  }\CommentTok{\# colour of the bars  }
                       \CommentTok{\# Hatches non{-}Python and R lighter}
\AttributeTok{density=}\FunctionTok{c}\NormalTok{(}\DecValTok{100}\NormalTok{, }\DecValTok{50}\NormalTok{, }\DecValTok{100}\NormalTok{, }\DecValTok{50}\NormalTok{, }\DecValTok{50}\NormalTok{, }\DecValTok{50}\NormalTok{, }\DecValTok{50}\NormalTok{, }\DecValTok{50}\NormalTok{, }\DecValTok{50}\NormalTok{, }\DecValTok{50}\NormalTok{)) }
\end{Highlighting}
\end{Shaded}

\pandocbounded{\includegraphics[keepaspectratio]{CS5702_W1_Lab_files/figure-latex/unnamed-chunk-3-1.pdf}}

\textbf{Exercise 2.4:} Hopefully you consider this bar chart to be an
improvement, but it's not perfect. Can you think of any further
improvements or clarifications? (Hint: some of the bars are quite hard
to read off the exact proportion.) Remember that if the goal of Data
Science is to \emph{communicate} findings, then how we present
information is of vital importance. For this reason we will return to
this topic in Week 4 to dig deeper on data visualisation.

\textbf{Extension Exercise 2.5:} The same
\href{https://www.onetonline.org/link/demand/15-2051.00}{article} also
gives the top 21 skills required and whether they are considered to be
``hot technology'', in an associated csv file. Can you think how you
might integrate this additional information in an effective
visualisation? Later in this worksheet we touch on how to read in files
to save scraping and typing \hyperref[WS_ReadingCSV]{here}.

\textbf{Extension exercise 2.6:} The above two plots are performed using
base R (which is perfectly adequate for simple charts). A powerful and
flexible alternative is the \{ggplot2\} package. How might you reproduce
these charts using \{ggplot2\}? NB This exercise is primarily intended
for those who already have some knowledge of R or who consider
themselves proficient programmers.

\section{3. Some Basic R practice}\label{WS3}

Essential tips for R:

\begin{itemize}
\tightlist
\item
  R is case-sensitive, so price and Price are different names.
\item
  In R everything that exists is an object.
\item
  Functions are also objects: mean refers to the function, while mean(x)
  calls it.
\item
  \texttt{\#} in front of a sentence or a word is treated as a comment
  and is not executed.
\item
  For help in R,
  \texttt{?\textless{}function\ or\ command\textgreater{}} e.g.,
  \texttt{?head}
\end{itemize}

\subsection{3.1 A simple variable assignment
example}\label{a-simple-variable-assignment-example}

In the previous task you may well have largely relied on the provided R
code. Now we move onto learning about basic R syntax and some coding
exercises for you to undertake.

Suppose I'm trying to calculate how much money I'm going to make from my
\emph{Modern Data} book. There are several different things I will need
to store:

\begin{itemize}
\tightlist
\item
  \texttt{sales} how many copies I'll sell\\
\item
  \texttt{price} the price per copy
\item
  \texttt{numberStudents} number of students in the class (assume I sell
  one copy per student)
\item
  \texttt{recommend} the proportion of students so impressed that they
  recommend the book to friends and family
\item
  \texttt{income} the income generated
\end{itemize}

Let's turn this into R.

\begin{Shaded}
\begin{Highlighting}[]
\CommentTok{\# Compute sales income for Martin\textquotesingle{}s book}

\NormalTok{price }\OtherTok{\textless{}{-}} \DecValTok{10}                                 \CommentTok{\# assume price is 10 UKP per book copy {-}}
                                            \CommentTok{\# actually it\textquotesingle{}s free but never mind!}
\NormalTok{numberStudents }\OtherTok{\textless{}{-}} \DecValTok{100}                       \CommentTok{\# number of students in the class}
\NormalTok{recommend }\OtherTok{\textless{}{-}} \FloatTok{0.5}                            \CommentTok{\# proportion of students who recommend }
                                            \CommentTok{\# the book}

\NormalTok{sales }\OtherTok{\textless{}{-}}\NormalTok{ numberStudents }\SpecialCharTok{*}\NormalTok{ (}\DecValTok{1} \SpecialCharTok{+}\NormalTok{ recommend)}
\NormalTok{income }\OtherTok{\textless{}{-}}\NormalTok{ sales }\SpecialCharTok{*}\NormalTok{ price}
\FunctionTok{paste}\NormalTok{(}\StringTok{"Martin will earn:"}\NormalTok{, income)          }\CommentTok{\# display a readable string with the paste function}
\end{Highlighting}
\end{Shaded}

\begin{verbatim}
## [1] "Martin will earn: 1500"
\end{verbatim}

Although we could just output \texttt{income}, I have added some text to
make the output of the code easier to interpret with the
\texttt{paste()} function to combine a character string literal
delimited by quotation marks and the calculated income.

\textbf{Exercise 3.1:} Change the number of students to 200 and the
recommendation rate to 0.2. What's my new estimated income?

\textbf{Exercise 3.2:} Update the sales income R code to handle tax.
Assume that income is taxed at a flat rate of 20\%. Modify the above R
code to output gross income and net (after tax) income. Using RStudio
you can edit this Week 1 Worksheet document.

\begin{Shaded}
\begin{Highlighting}[]
\CommentTok{\# Your extended code needs to go here }
\end{Highlighting}
\end{Shaded}

HINT: I recommend you define a new numeric variable to store the tax
rate and another for net income. For clarity you might want to rename
\texttt{income} to make it obvious that it stores gross income.

\textbf{Extension Exercise 3.3:} How would you answer the question: it's
possible to write a much shorter R program without any extra variables,
so why bother?

It's very \emph{important} to get into the habit of writing easy to
follow R code. For this reason I have added comments and tried to select
meaningful variable names.

\subsection{3.2 Using vectors}\label{using-vectors}

Recall that a Vector is a repeated group of the \textbf{same} data type,
e.g., a vector of character strings.

\begin{Shaded}
\begin{Highlighting}[]
\CommentTok{\# Define a new vector}
\NormalTok{greetings }\OtherTok{\textless{}{-}} \FunctionTok{c}\NormalTok{(}\StringTok{"Hi!"}\NormalTok{, }\StringTok{"Hello"}\NormalTok{, }\StringTok{"Good morning"}\NormalTok{, }\StringTok{"Saludo"}\NormalTok{, }\StringTok{"Hej"}\NormalTok{, }\StringTok{"¡Hola!"}\NormalTok{)}
\NormalTok{greetings}
\end{Highlighting}
\end{Shaded}

\begin{verbatim}
## [1] "Hi!"          "Hello"        "Good morning" "Saludo"       "Hej"         
## [6] "¡Hola!"
\end{verbatim}

We can also check whether the type of \texttt{greetings} is a vector.

\begin{Shaded}
\begin{Highlighting}[]
\CommentTok{\# Check greetings is a vector}
\FunctionTok{is.vector}\NormalTok{(greetings)}
\end{Highlighting}
\end{Shaded}

\begin{verbatim}
## [1] TRUE
\end{verbatim}

\textbf{Exercise 3.4:} Check the type of the variable \texttt{price}
(from Section 3.1). What output do you expect?

To fully exploit the power of vectors we can access subsets of elements.
To do this we use the square bracket notation
\texttt{{[}\textless{}n\textgreater{}{]}} where \(n\) is the position of
the required vector element.

\begin{Shaded}
\begin{Highlighting}[]
\CommentTok{\# Different ways to index vector elements}
\NormalTok{greetings[}\DecValTok{2}\NormalTok{]        }\CommentTok{\# Access the 2nd element}
\end{Highlighting}
\end{Shaded}

\begin{verbatim}
## [1] "Hello"
\end{verbatim}

\begin{Shaded}
\begin{Highlighting}[]
\NormalTok{n }\OtherTok{\textless{}{-}} \DecValTok{3}
\NormalTok{greetings[n]        }\CommentTok{\# You can use a variable as an index }
\end{Highlighting}
\end{Shaded}

\begin{verbatim}
## [1] "Good morning"
\end{verbatim}

\begin{Shaded}
\begin{Highlighting}[]
\NormalTok{greetings[}\DecValTok{1}\SpecialCharTok{:}\DecValTok{3}\NormalTok{]      }\CommentTok{\# Or you can specify a range of vector elements}
\end{Highlighting}
\end{Shaded}

\begin{verbatim}
## [1] "Hi!"          "Hello"        "Good morning"
\end{verbatim}

Notice in the above R chunk, that we can access more than one element at
a time by using the \texttt{:} operator to return a range of elements.
The range can be specified by literals as in \texttt{{[}1:3{]}} or
integer\footnote{If the index variables are not integers R will do its
  best to coerce the values to integers.} variables e.g.,
\texttt{{[}m:n{]}}.

\subsection{3.3 Using data frames}\label{using-data-frames}

Data frames are amongst the most useful of the different data
structures. These are two dimensional and unlike 2-dimensional matrices
can comprise different constituent data types.

If you're not very confident about data frames I suggest you view my
short(ish)
\href{https://www.loom.com/share/f57102f824e948d2b357f1909ad070a1}{video
here}.

Let's start by examining a built-in data set \texttt{mtcars} that stores
information about motor cars from the 1970s that is stored as a data
frame.

\begin{Shaded}
\begin{Highlighting}[]
\CommentTok{\# Assign the built in data object mtcars to a new data frame df}
\NormalTok{df }\OtherTok{\textless{}{-}}\NormalTok{ mtcars}

\CommentTok{\# View the first 6 rows of the data frame df}
\FunctionTok{head}\NormalTok{(df)}
\end{Highlighting}
\end{Shaded}

\begin{verbatim}
##                    mpg cyl disp  hp drat    wt  qsec vs am gear carb
## Mazda RX4         21.0   6  160 110 3.90 2.620 16.46  0  1    4    4
## Mazda RX4 Wag     21.0   6  160 110 3.90 2.875 17.02  0  1    4    4
## Datsun 710        22.8   4  108  93 3.85 2.320 18.61  1  1    4    1
## Hornet 4 Drive    21.4   6  258 110 3.08 3.215 19.44  1  0    3    1
## Hornet Sportabout 18.7   8  360 175 3.15 3.440 17.02  0  0    3    2
## Valiant           18.1   6  225 105 2.76 3.460 20.22  1  0    3    1
\end{verbatim}

\textbf{Exercise 3.5:} How many variables and observations does
\texttt{df} have? HINT: There are functions \texttt{nrow()} and
\texttt{ncol()} if you can't be bothered to count.

You can use the \texttt{\$} operator to specify a column/variable from a
data frame, for instance:

\begin{Shaded}
\begin{Highlighting}[]
\CommentTok{\# Display the observations for cylinder numbers from df}
\NormalTok{df}\SpecialCharTok{$}\NormalTok{cyl}
\end{Highlighting}
\end{Shaded}

\begin{verbatim}
##  [1] 6 6 4 6 8 6 8 4 4 6 6 8 8 8 8 8 8 4 4 4 4 8 8 8 8 4 4 4 8 6 8 4
\end{verbatim}

\begin{Shaded}
\begin{Highlighting}[]
\CommentTok{\# Apply the median function}
\FunctionTok{median}\NormalTok{(df}\SpecialCharTok{$}\NormalTok{cyl)}
\end{Highlighting}
\end{Shaded}

\begin{verbatim}
## [1] 6
\end{verbatim}

\textbf{Exercise 3.6:} Using the \texttt{\$} operator determine the mean
and minimum mpg for cars in df.

\textbf{Exercise 3.7:} Update the number of gears for the Mazda RX4 from
4 to 6. HINT: \texttt{df\$gear} is a vector so you can use an index to
reference a particular element.

So we can use the dollar operator to identify columns/variables. What
about rows/observations? There are several ways to achieve this. Using
indices via the \texttt{{[}\ {]}} notation notice that since a data
frame has two dimensions two arguments are needed, but since we want the
entire row we can leave the argument after the comma blank.

\begin{Shaded}
\begin{Highlighting}[]
\CommentTok{\# To retrieve the entire 2nd row we could have.}
\NormalTok{df[}\DecValTok{2}\NormalTok{,]}
\end{Highlighting}
\end{Shaded}

\begin{verbatim}
##               mpg cyl disp  hp drat    wt  qsec vs am gear carb
## Mazda RX4 Wag  21   6  160 110  3.9 2.875 17.02  0  1    4    4
\end{verbatim}

\textbf{Exercise 3.8:} Modify the above code to retrieve just columns
1-3 of row 2. HINT: remember the \texttt{:} operators from when we were
manipulating vectors.

The other way to retrieve specific observations is when the rows have
names e.g., ``Mazda RX4'', or in this case:

\begin{Shaded}
\begin{Highlighting}[]
\CommentTok{\# Retrieve the observation named "Valiant"}
\NormalTok{df[}\StringTok{"Valiant"}\NormalTok{,]}
\end{Highlighting}
\end{Shaded}

\begin{verbatim}
##          mpg cyl disp  hp drat   wt  qsec vs am gear carb
## Valiant 18.1   6  225 105 2.76 3.46 20.22  1  0    3    1
\end{verbatim}

\textbf{Exercise 3.9:} Change the number of cylinders for the Valiant to
12.

\textbf{Quiz:} Please check your understanding of this chapter via this
\href{https://forms.gle/m6J1ka6LECvaDKGu5}{quick quiz} (10 simple
multiple choice questions with answers).

\subsection{3.4 Reading in a csv file}\label{WS_ReadingCSV}

So far we have made use of data that is already in R or input our own.
In practice often data will be provided in .csv format and R makes it
easy to read it in.

\begin{itemize}
\tightlist
\item
  To make your life easier we have hosted the worms file on a GitHub
  repository.\\
\item
  Run the code below to read data from the hard coded url.
\end{itemize}

\begin{Shaded}
\begin{Highlighting}[]
\CommentTok{\#}
\NormalTok{worms }\OtherTok{\textless{}{-}} \FunctionTok{read.csv}\NormalTok{(}\StringTok{"https://raw.githubusercontent.com/mjshepperd/CS5702{-}Data/master/worms.csv"}\NormalTok{)}
\end{Highlighting}
\end{Shaded}

If you want to check what the data looks like:

\begin{Shaded}
\begin{Highlighting}[]
\FunctionTok{head}\NormalTok{(worms,}\AttributeTok{n=}\DecValTok{10}\NormalTok{)}
\end{Highlighting}
\end{Shaded}

\begin{verbatim}
##         Field.Name Area Slope Vegetation Soil.pH  Damp Worm.density
## 1      Nashs.Field  3.6    11  Grassland     4.1 FALSE            4
## 2   Silwood.Bottom  5.1     2     Arable     5.2 FALSE            7
## 3    Nursery.Field  2.8     3  Grassland     4.3 FALSE            2
## 4      Rush.Meadow  2.4     5     Meadow     4.9  TRUE            5
## 5  Gunness.Thicket  3.8     0      Scrub     4.2 FALSE            6
## 6         Oak.Mead  3.1     2  Grassland     3.9 FALSE            2
## 7     Church.Field  3.5     3  Grassland     4.2 FALSE            3
## 8          Ashurst  2.1     0     Arable     4.8 FALSE            4
## 9      The.Orchard  1.9     0    Orchard     5.7 FALSE            9
## 10   Rookery.Slope  1.5     4  Grassland     5.0  TRUE            7
\end{verbatim}

In QDA you will see how to do some more descriptive data analysis.

\subsection{Extension Appendix 1: extra notes on
YAML}\label{YAML_appendix}

YAML rhymes with camel! It is an acronym for ``\textbf{Y}AML
\textbf{A}in't \textbf{M}arkup \textbf{L}anguage''.

An R Markdown file has three components:

\begin{enumerate}
\def\labelenumi{\arabic{enumi}.}
\tightlist
\item
  Header or front matter written in YAML
\item
  Markdown: this is text to describe your R analysis
\item
  Code chunks: Chunks of embedded R code that can be executed and
  rendered
\end{enumerate}

An RMarkdown file \emph{must} start with a header written using YAML
syntax. There are four default elements in the RStudio YAML header:

\begin{verbatim}
title: The title of your document. NB this is not the same as the file name.
author: Who wrote the document.
date: By default this is the date that the file is created.
output: What format will the output be in.  We will generally use html.  Other options are pdf, various slide formats, etc.
\end{verbatim}

The YAML is delineated by \texttt{-\/-\/-} and \texttt{-\/-\/-}.

However, if you're unsure, then your best strategy is to copy and adapt,
e.g., change the date, title or author.

\subsection{Extension Appendix 2: Pie charts with
ggplot2}\label{ggplot_appendix}

Remember that many data scientists have strong reservations about pie
charts. In this case a bar plot is probably more effective.

NB This optional example will run only if \{ggplot2\} is installed. If
no chart appears, run \texttt{install.packages("ggplot2")} in the
Console and knit the document again.

\begin{Shaded}
\begin{Highlighting}[]
\CommentTok{\# Reproduce the data from the 2025 required tools for Data Scientists}
\CommentTok{\# using ggplot2}

\CommentTok{\# load the ggplot2 library}
\FunctionTok{library}\NormalTok{(ggplot2)}

\NormalTok{skill }\OtherTok{\textless{}{-}} \FunctionTok{c}\NormalTok{(}
  \StringTok{"Python"}\NormalTok{, }\StringTok{"SQL"}\NormalTok{, }\StringTok{"R"}\NormalTok{, }\StringTok{"Tableau"}\NormalTok{, }\StringTok{"Power BI"}\NormalTok{,}
  \StringTok{"AWS"}\NormalTok{, }\StringTok{"Azure"}\NormalTok{, }\StringTok{"TensorFlow"}\NormalTok{, }\StringTok{"PyTorch"}\NormalTok{, }\StringTok{"SAS"}
\NormalTok{)}
\NormalTok{percent }\OtherTok{\textless{}{-}} \FunctionTok{c}\NormalTok{(}\DecValTok{66}\NormalTok{, }\DecValTok{51}\NormalTok{, }\DecValTok{34}\NormalTok{, }\DecValTok{22}\NormalTok{, }\DecValTok{19}\NormalTok{, }\DecValTok{17}\NormalTok{, }\DecValTok{13}\NormalTok{, }\DecValTok{11}\NormalTok{, }\DecValTok{10}\NormalTok{, }\DecValTok{9}\NormalTok{)}

\CommentTok{\# Merge into a single dataframe which makes ggplot easier}
\NormalTok{dsTools }\OtherTok{\textless{}{-}} \FunctionTok{data.frame}\NormalTok{(skill,percent)}

\CommentTok{\# Generate the pie chart}
\FunctionTok{ggplot}\NormalTok{(dsTools, }\FunctionTok{aes}\NormalTok{(}\AttributeTok{x=}\StringTok{""}\NormalTok{, }\AttributeTok{y=}\NormalTok{percent, }\AttributeTok{fill=}\NormalTok{skill)) }\SpecialCharTok{+}
  \FunctionTok{geom\_bar}\NormalTok{(}\AttributeTok{stat=}\StringTok{"identity"}\NormalTok{, }\AttributeTok{width=}\DecValTok{1}\NormalTok{) }\SpecialCharTok{+}
  \FunctionTok{coord\_polar}\NormalTok{(}\StringTok{"y"}\NormalTok{, }\AttributeTok{start=}\DecValTok{0}\NormalTok{) }\SpecialCharTok{+}
  \FunctionTok{theme\_void}\NormalTok{() }\SpecialCharTok{+}                 \CommentTok{\# remove background, grid, numeric labels}
  \FunctionTok{labs}\NormalTok{(}\AttributeTok{title =} \StringTok{"In demand data scientist tools (2025)"}\NormalTok{) }\SpecialCharTok{+}
  \FunctionTok{theme}\NormalTok{(}\AttributeTok{plot.title =} \FunctionTok{element\_text}\NormalTok{(}\AttributeTok{size =} \DecValTok{10}\NormalTok{))}
\end{Highlighting}
\end{Shaded}

\pandocbounded{\includegraphics[keepaspectratio]{CS5702_W1_Lab_files/figure-latex/unnamed-chunk-15-1.pdf}}

\textbf{Extension Exercise A2} Rewrite the above ggplot2 to produce a
bar chart. Add suitable labels.

\end{document}
